%----------------------------------------------------------------------
% 中文摘要
%----------------------------------------------------------------------

% 把中文摘要寫在裡面
\begin{zhAbstract}

邊緣電漿子（Edge Plasmons）技術能將光場局域化於幾何邊緣，並提供極高的微觀動量與場增強效應，在低維度奈米光學與二維材料物理中具備應用潛力。然而，傳統表面電漿極化子（SPPs）受限於動量匹配範圍，難以有效激發或點亮單層過渡金屬二硫族化合物（Monolayer Transition Metal Dichalcogenides, TMDs）中具有高動量點（High-momentum points）的特定激子（Excitons）態。

本研究結合時域有限差分法（Finite-Difference Time-Domain, FDTD），利用 MEEP 模擬平台，分析不同金屬薄膜厚度對邊緣電漿子模式之幾何色散關係的調控效應，藉此探討其與 TMDs 能帶結構之動量匹配機制。

模擬結果顯示，改變金屬薄膜厚度能有效調控邊緣電漿子的色散軌跡。此電漿子色散曲線與 TMDs 激子色散關係之交點，即為兩者滿足能量與動量守恆、發生高效光物質交互作用之共振點。透過此幾何調控機制，本研究證實在交點共振區域能提升高動量激子之光學躍遷速率（Transition rate），成功實現激子點亮效應。此研究結果展現了幾何邊緣結構在微觀動量工程中的應用潛力，並可作為開發二維材料基光電與量子元件的有效手段。

    % 這個Command會自動幫你把Config裡面設定的東東填進來
    \zhAbsKeywords
\end{zhAbstract}
